% ============================================================
% Abstract — OR journal style: single paragraph + keywords
% ============================================================
\begin{abstract}
Diffusion models generate samples by traversing a denoising trajectory,
a sequence of noise-reduction steps that progressively transforms random noise
into data of the target distribution.
Once diffusion models are pre-trained and deployed into inference, i.e.\ using
compute to run the model and generate data samples, a practical question is how
to use a given compute budget to generate higher quality data samples.
Standard approaches allocate compute uniformly across the denoising steps,
but this is wasteful: early steps that determine global structure benefit more
from additional compute than late steps that refine local details.
We formalize this as a budget allocation problem: given $B$ total function
evaluations across $T$ denoising steps, find the allocation that
maximizes sample quality.
We prove that the optimal allocation follows
a water-filling rule that concentrates compute on steps where additional
refinement yields the largest quality gains.
On Stable Diffusion, EDM, and flow-based models, our method achieves
20--50\% budget reduction while maintaining the same quality as
uniform-allocation benchmarks, and consistently exceeds them at equal budget.
\end{abstract}

\noindent\textit{Key words}: diffusion models; inference-time scaling; computational budget allocation; noise trajectory search; global scheduling; resource allocation; sequential decision making

\vspace{1em}
